\subsection{Subsistema de Alimentação e Regulação}
Este subsistema gerencia a conversão e distribuição de energia para garantir que cada componente opere dentro de sua faixa de tensão nominal, prevenindo danos ao hardware.

\begin{itemize}
    \item \textbf{Entrada de Energia}: O sistema recebe alimentação de 5V via porta Micro-USB, conectada diretamente ao computador ou fonte externa.
    \item \textbf{Regulação de Tensão (LDO)}: A placa DevKit utiliza um regulador de baixa queda (\textit{Low-Dropout} - LDO) para converter os 5V da entrada em 3,3V estáveis, necessários para o funcionamento do chip ESP32 e do sensor AHT20.
    \item \textbf{Barramento de 5V}: Este barramento alimenta o módulo de cartão SD, que por sua vez possui regulação própria e \textit{level shifters} internos para garantir a integridade dos sinais lógicos.
\end{itemize}

\subsection{Subsistema de Sensoriamento e Armazenamento}
Responsável por ler as informações do ambiente e salvá-las fisicamente.

\begin{itemize}
    \item \textbf{Sensor}: O AHT20 mede temperatura e umidade usando o protocolo I2C.
    \item \textbf{Gravação}: As leituras são salvas no cartão SD em arquivos \texttt{.csv} via protocolo SPI, garantindo que os dados fiquem guardados mesmo sem internet.
\end{itemize}

\subsection{Subsistema de Interface e Atuação}
Proporciona ao usuário uma visualização imediata do status operacional do dispositivo através de sinalizações luminosas.

\begin{itemize}
    \item \textbf{Sinalização Visual}: Três LEDs (Verde, Amarelo e Vermelho) compõem a interface física, permitindo identificar visualmente os estados de operação do hardware, estes leds funcionarão como: VERDE (ESTADO LIGADO), AMARELO (ESTADO AUTOMÁTICO), VERMELHO (ESTADO DESLIGADO)
    \item \textbf{Proteção Física}: Resistores limitadores de corrente são empregados em série com os LEDs, protegendo as saídas digitais do microcontrolador contra drenagem excessiva de corrente.
\end{itemize}

\subsection{Subsistema de Comunicação e Arquitetura de Software}
Este bloco gerencia a conectividade sem fio e a entrega da interface de usuário, permitindo o monitoramento remoto e a configuração do dispositivo em tempo real.

\begin{itemize}
    \item \textbf{Conectividade Wi-Fi}: O microcontrolador opera nos modos \textit{Station} (conectado a uma rede local) ou \textit{Access Point} (criando sua própria rede para configuração inicial), utilizando a pilha TCP/IP nativa do ESP32.
    \item \textbf{Servidor Web Assíncrono}: Utiliza-se a biblioteca \texttt{ESPAsyncWebServer} para processar requisições HTTP de forma não bloqueante. Isso garante que a entrega das páginas do site não interrompa as tarefas críticas, como a leitura do sensor AHT20 ou a escrita de logs no cartão SD.
    \item \textbf{Gerenciamento de Arquivos (LittleFS)}: Os arquivos de \textit{front-end} (HTML, CSS e JavaScript) são armazenados em uma partição da memória flash do microcontrolador organizada pelo sistema de arquivos LittleFS. Essa abordagem separa a lógica de controle (C++) da interface visual, otimizando o uso da memória RAM.
    \item \textbf{Interatividade (JSON e AJAX)}: A troca de dados entre o hardware e o navegador ocorre via pacotes JSON. Através de requisições assíncronas (AJAX), o site atualiza os valores de temperatura e umidade na tela automaticamente, eliminando a necessidade de recarregar a página pelo usuário.
\end{itemize}
